\chapter{Couche de persistance}\label{chap:persistance}

\section{Les entités et leurs correspondances}
Huit classes sont persistantes (\code{Internaute}, \code{Client}, \code{Administrateur}, \code{Vitrine},
\code{Categorie}, \code{Produit}, \code{Panier}, \code{LignePanier}). Le tableau~\ref{tab:mapping} récapitule les
décisions de mapping ; le \textit{pourquoi} de chacune est donné ensuite.

\begin{table}[htbp]
\centering
\caption{Décisions de mapping}\label{tab:mapping}
\begin{tabularx}{\linewidth}{@{}lXX@{}}
\toprule
\textbf{Sujet} & \textbf{Décision} & \textbf{Raison}\\
\midrule
Accès aux données & Sur les champs (\code{@Id} sur l'attribut) & Les accesseurs peuvent contenir de la logique (normalisation de l'e-mail) sans être appelés par JPA.\\
Montants & \code{BigDecimal}, échelle 2 & Pas d'erreur d'arrondi, contrairement à \code{double}.\\
Concurrence & \code{@Version} sur \code{Produit} & Détecte deux modifications simultanées sans verrou en base.\\
Client -- Panier & \code{@OneToOne}, cascade \code{ALL}, \code{orphanRemoval} & Le panier naît et meurt avec son client.\\
Panier -- Lignes & \code{@OneToMany} avec \code{@CascadeOnDelete} & Supprimer un panier supprime ses lignes, y compris au niveau SQL (\code{ON DELETE CASCADE}).\\
Ligne -- Produit & \code{@ManyToOne}, sans cascade & Une ligne référence un produit, elle n'en est pas propriétaire.\\
Listes exposées & Copies en lecture seule & Aucun code externe ne modifie une collection sans passer par les règles métier.\\
\bottomrule
\end{tabularx}
\end{table}

\textbf{Le verrouillage optimiste.} Deux administrateurs peuvent ouvrir le même produit, ou un client peut
ajouter au panier le dernier exemplaire pendant qu'on modifie son stock. Un verrou pessimiste bloquerait des lignes
en base pendant la réflexion d'un utilisateur ; ce n'est pas acceptable dans une application web. Avec
\code{@Version}, chaque mise à jour incrémente un numéro, et une écriture fondée sur un numéro périmé échoue. L'application
traduit cet échec en message clair (« Cet élément a été modifié par quelqu'un d'autre... »). Dans les opérations
de panier, le produit est lu avec \code{LockModeType.OPTIMISTIC} : si son numéro change avant la validation, la
transaction échoue au lieu de laisser passer une quantité calculée sur un stock périmé.

\textbf{Le prix figé.} Une ligne de panier recopie le prix du produit au moment de l'ajout. Si l'administrateur
change le prix ensuite, le panier du client ne bouge pas, ce qui évite qu'un total change sous ses yeux.

\section{L'unité de persistance}
Le fichier \code{persistence.xml} ne contient ni mot de passe ni adresse de base. Il décrit seulement ce qui ne varie
pas d'une machine à l'autre.

\lstinputlisting[language=XML,caption={persistence.xml (sans les commentaires)},basicstyle=\ttfamily\scriptsize]{extraits/persistence.txt}

Les choix notables sont les suivants.
\begin{itemize}
  \item \code{RESOURCE\_LOCAL} : Tomcat n'est pas un serveur JTA, l'application gère donc ses transactions
        elle-même (section~\ref{sec:transactions}).
  \item Liste explicite des classes : l'analyse automatique du classpath n'est pas fiable sous Tomcat.
  \item \code{shared-cache-mode NONE} : pas de cache partagé, donc ce que l'on voit dans MySQL est exactement ce que
        l'application lit. C'est un choix de simplicité, adapté à un projet pédagogique.
  \item \code{eclipselink.weaving=false} : il n'y a pas d'agent de tissage sous Tomcat. Conséquence connue,
        signalée par un avertissement au démarrage : les relations \code{@OneToOne} ne sont pas chargées
        paresseusement.
\end{itemize}

\textbf{Identifiants hors du dépôt.} L'adresse de la base, l'utilisateur et le mot de passe viennent d'un fichier
\code{db.properties}, ignoré par Git. Le dépôt ne contient qu'un modèle, \code{db.properties.example}. Une fuite d'identifiants
dans un historique Git est difficile à effacer : la précaution est donc prise dès le premier commit. La classe
\code{PersistenceConfig} lit ce fichier et construit les propriétés de la fabrique d'\code{EntityManager}.

\section{Création du schéma et données de test}
Le schéma n'est pas laissé au hasard du démarrage de l'application. Un outil en ligne de commande, \code{SchemaTool},
le crée explicitement :
\begin{enumerate}
  \item il demande de taper \code{RESET} avant de supprimer et recréer les tables (option \code{--yes} pour les scripts) ;
  \item il applique les contraintes \code{CHECK} de \code{constraints.sql}, que JPA ne sait pas exprimer ;
  \item il charge les données de démonstration : une vitrine, quatre catégories, douze produits dont un en rupture
        (\code{MAI-002}), un à trois exemplaires (\code{INF-003}) et un retiré de la vente (\code{MOD-003}), un administrateur
        et deux clients.
\end{enumerate}
Séparer la création du schéma du démarrage de l'application évite qu'un redéploiement modifie la base. Pour la même
raison, \code{jpa.ddl} vaut \code{none} lors de l'exécution normale.

\fig{01-persistence-xml}{Le fichier persistence.xml ouvert dans Eclipse (onglet Source)}{fig:persistencexml}
\fig{02-schema-tool-console}{Sortie de SchemaTool : contraintes appliquées, données chargées, bilan}{fig:schematool}

\section{Vérification dans MySQL}
La base a été inspectée directement, sans passer par l'application, pour valider le mapping avant d'écrire la
moindre servlet.

\begin{itemize}
  \item Les six tables attendues existent (figure~\ref{fig:showtables}) : \code{categorie}, \code{internaute},
        \code{ligne\_panier}, \code{panier}, \code{produit}, \code{vitrine}.
  \item \code{produit.prix} est un \code{decimal(10,2)}, \code{produit.version} un \code{bigint}, et
        \code{reference} porte une contrainte d'unicité (figure~\ref{fig:describe}).
  \item Dans \code{internaute}, les colonnes propres aux sous-classes (\code{adresse}, \code{telephone},
        \code{matricule}, \code{panier\_id}) sont nullables : c'est la contrepartie de \code{SINGLE\_TABLE}.
  \item \code{ligne\_panier} porte \code{ON DELETE CASCADE}, l'unicité (\code{panier\_id}, \code{produit\_id}) et deux
        \code{CHECK} (figure~\ref{fig:showcreate}).
  \item Les accents sont intacts (« mécanique », « Écharpe »), ce qui valide la chaîne UTF-8 complète, de la
        base jusqu'à l'affichage (figures~\ref{fig:internautes} et~\ref{fig:produits}).
\end{itemize}

\fig{03-show-tables}{Les six tables créées}{fig:showtables}
\fig{04-describe-tables}{Structure des tables produit et internaute}{fig:describe}
\fig{05-show-create-ligne-panier}{Définition de ligne\_panier : clé unique, cascade et contraintes CHECK}{fig:showcreate}
\fig{06-donnees-internautes}{Les internautes de démonstration et leur rôle}{fig:internautes}
\fig{07-donnees-produits}{Les douze produits de démonstration}{fig:produits}

\section{Les contraintes refusent bien les données incorrectes}
Trois écritures invalides ont été tentées à la main. La base les refuse toutes, ce qui montre qu'elle protège ses
données même si une règle était oubliée dans le code Java (tableau~\ref{tab:refus}).

\begin{table}[htbp]
\centering
\caption{Écritures refusées par la base}\label{tab:refus}
\begin{tabularx}{\linewidth}{@{}Xll@{}}
\toprule
\textbf{Tentative} & \textbf{Code MySQL} & \textbf{Contrainte}\\
\midrule
Mettre un stock à $-1$ & 3819 & \code{chk\_produit\_stock}\\
Insérer un second internaute avec la même adresse & 1062 & \code{uk\_internaute\_email}\\
Supprimer un produit présent dans un panier & 1451 & clé étrangère de \code{ligne\_panier}\\
\bottomrule
\end{tabularx}
\end{table}

\fig{08-contraintes-refusees}{Les trois refus et leurs codes d'erreur}{fig:refus}

\textbf{Une observation corrigée.} La première vérification a montré que \code{type\_internaute}, la colonne
discriminante, était nullable : EclipseLink la crée ainsi. Un insertion manuelle pouvait donc créer un utilisateur
sans rôle. Le fichier \code{constraints.sql} impose désormais \code{NOT NULL} et un \code{CHECK} limitant les valeurs à
\code{CLIENT} et \code{ADMIN} (règle R11).
